\section{An unramified Gaussian cover and an exact Jacobian torsion class}
\label{ug-section}

Let $g\geq3$ be an odd integer, and write $D=D_{g,1}$ and
$C=C_{g,1}$ for the curves defined over $\mathbb Q$ in
Section~\ref{rl-section}. Forgetting the Gaussian coordinate gives
$\pi:C\to D$. Work also over the fixed field
$K=\mathbb Q(i,\sqrt{-3})$. With the four quadratics of
Lemma~\ref{bk-branches}, the function fields are
\begin{equation}\label{ug-fields}
 K(D)=K(t,Y),\quad Y^g=\frac{B_1B_2}{B_3B_4},\qquad
 K(C)=K(t,Y,Z),\quad Z^g=\frac{B_1B_4}{B_2B_3}.
\end{equation}
The eight simple roots of the four $B_j$ are distinct and finite.
The preceding geometric degree computations give degrees $g$ and
$g^2$ over $\mathbb P^1_t$.

\begin{theorem}[An everywhere unramified relative cover]
\label{ug-etale}
The map $\pi$ is finite \'etale of degree $g$ over $\mathbb Q$, and
\begin{equation}\label{ug-genera}
 g(D)=3g-3,\qquad g(C)=1+3g^2-4g=1+g\bigl(g(D)-1\bigr).
\end{equation}
\end{theorem}
\begin{proof}
The function-field inclusion extends to a finite nonconstant map
of smooth projective curves, of degree $g$.
On the base $\mathbb P^1_t$, the Kummer function defining $D$
has a simple zero at roots of $B_1,B_2$
and a simple pole at roots of $B_3,B_4$. Over an algebraic closure,
each has a unique point above it, of ramification index $g$.
At other points, including infinity, its valuation is zero and the
local unit Kummer equation is unramified. Riemann--Hurwitz gives
$2g(D)-2=-2g+8(g-1)=6g-8$.

On $C$ the local valuation pairs at those roots are
$(1,1),(1,-1),(-1,-1),(-1,1)$. Extracting roots of local units over
an algebraically closed field shows that both Kummer extensions
lie in one extension generated by a $g$-th root of a uniformizer.
Its index is exactly $g$. Elsewhere the index is one.
Indices multiply in the tower $C\to D\to\mathbb P^1$, so the
relative index is one everywhere. In characteristic zero this
finite map of smooth curves is \'etale; this property descends to
$\mathbb Q$. The genus of $C$ is also
Theorem~\ref{bk-degree-genus}, and agrees with the unramified
Riemann--Hurwitz formula \cite{StacksCurves2026}.
\end{proof}

\begin{theorem}[One Kummer function and an exact torsion divisor]
\label{ug-exact-torsion}
Put $f=B_4/B_2$ and $m=(g+1)/2$. Then
\begin{equation}\label{ug-single-kummer}
 W=\frac{(Z/Y)^m}{f},\qquad W^g=f,\qquad W^2=Z/Y,\qquad
 K(C)=K(D)(W),\quad Z=YW^2.
\end{equation}
For $j=2,4$, let $R_j$ be the sum of the unique points of $D$
above the two roots of $B_j$. These are degree-two divisors
defined over $K$, and
\[
 \Theta=[R_4-R_2]\in\operatorname{Pic}^0(D_K)
\]
has exact order $g$, also geometrically. Its natural image in
$\operatorname{Jac}(D)(K)$ has exact order $g$.
\end{theorem}
\begin{proof}
The factorizations give
\[
 \frac{G/\bar G}{E/\bar E}=(B_4/B_2)^2=f^2.
\]
Thus $(Z/Y)^g=f^2$, and $2m=g+1$ proves both identities for $W$
in \eqref{ug-single-kummer}. Conversely $Z=YW^2$ gives the original
second equation, proving equality of the fields over $K$ itself.

The roots of $B_j$ form a Galois-stable pair and each has a unique
point above it, so $R_j$ is defined over $K$. Pullback multiplies
the zero and pole valuations of $f$ by $g$, and infinity has
valuation zero. Therefore
\begin{equation}\label{ug-divisor}
 \operatorname{div}_D(f)=g(R_4-R_2).
\end{equation}
Suppose over an algebraic closure $k$ the class had order $d<g$.
Then $d\mid g$ and some $h\in k(D)^*$ would have
$\operatorname{div}(h)=d(R_4-R_2)$. Since $D$ is projective and
geometrically integral, $f/h^{g/d}$ is a nonzero constant.
Taking a constant root in $k$, write $f=(ch)^{g/d}$.
If $r=g/d>1$, then
\[
 T^g-f=\prod_{\xi\in\mu_r(k)}(T^d-\xi ch).
\]
Every root has degree at most $d$ over $k(D)$, contradicting the
geometric degree $g$ of $k(C)/k(D)$ from Theorem~\ref{ug-etale}.
This proves exact geometric order, and \eqref{ug-divisor} gives
exact order over $K$ too.

The divisor gives an actual degree-zero line bundle over $K$,
hence a $K$-rational point under the natural Picard-to-Jacobian
map \cite[Theorem~1.1 and Remark~1.5]{MilneJV2021}.
After base change to $k$ this map is the usual identification
with the Jacobian. Exact geometric order therefore proves the
claimed order of the $K$-rational image. No rational base point
on $D$, or equality of all of
$\operatorname{Pic}^0(D_K)$ with $\operatorname{Jac}(D)(K)$,
is assumed.
\end{proof}

There is a precise arithmetic interpretation. Locally,
\eqref{ug-divisor} writes $f=s^g u$ with $s$ a rational local
equation for $R_4-R_2$ and $u$ a regular unit. The normalization
adjoins $W/s$ satisfying $(W/s)^g=u$. This is finite \'etale,
and scaling $W$ gives a torsor under the group scheme $\mu_g$
over $D_K$. Only geometrically do all $g$ scalar deck
transformations occur. We have not assumed $\mu_g\subset K$,
nor asserted a constant cyclic arithmetic deck group of order $g$.

\begin{corollary}[The actual positive rational fibers]\label{ug-positive}
Write $D(\mathbb Q)_{>1}$ and $C(\mathbb Q)_{>1}$ for points with
finite rational base coordinate $t>1$. The map
\[
 \pi:C(\mathbb Q)_{>1}\longrightarrow D(\mathbb Q)_{>1}
\]
is a bijection. If $\gcd(g,6)=1$, both sets are in bijection with
the ordered positive primitive seeds
$a^2+ab+b^2=U^2$, $F(a,b)=Q^g$, with $U,Q\geq7$.
\end{corollary}
\begin{proof}
This is exactly Theorem~\ref{rl-unique-lift}. Existence comes from
the actual Gaussian decomposition, while uniqueness uses
$\mu_g(K_G)=\{1\}$ for odd $g$. The geometric degree is not the
number of rational lifts. The additional coprimality condition is
the one needed to absorb the Eisenstein unit and obtain the
single-curve ordered-seed bijection.
\end{proof}

Thus the nontrivial \'etale cover and its exact-order Jacobian
class remove no point of this positive rational locus. Every
such fiber has a rational lift, and after base change to $K$
its $\mu_g$-torsor has a $K$-point. The class provides an explicit
invariant for a fixed-exponent descent or Jacobian computation.
Its ambient dimension is $3g-3$ and grows with $g$; no
fixed-dimension uniform torsion conclusion applies.
No point-height upper bound, full rational-point computation,
odd-exponent emptiness or ABC theorem is obtained. Ramified
first norms and moving nonunit residuals remain separate branches.
